\BookRecto
\begingroup
\renewcommand{\runningtitle}{Pronunciation key}
\hyphenpenalty=10000
\exhyphenpenalty=10000
\emergencystretch=3em
\newcommand{\reserve}[1]{\par\penalty-100\begingroup\dimen0=\pagegoal\advance\dimen0 by-\pagetotal\ifdim\dimen0<#1\newpage\fi\endgroup}
\newcommand{\itempara}[2]{\par\addvspace{3pt}\begingroup\raggedright\noindent\hangindent=1.8em\hangafter=1\makebox[1.8em][l]{#1.}#2\par\endgroup}
% book-source-body:begin

\addfontfeatures{Ligatures=TeXOff}
\raggedright
\booktitle{\BookComponent{pronunciation-key}{Pronunciation practice: key}SEPARATE KEY}{Pronunciation practice: key}

% mdparagraph:P001:begin
These answers follow the guide's chosen reading convention. English cues identify particular sounds, not substitute spellings for whole Old English words. The permitted American \textbf{r} and practical short-vowel approximations are acceptable. Meanings and grammatical forms are supplied throughout; the tasks assess pronunciation.
% mdparagraph:P001:end

\reserve{4\baselineskip}
% mdparagraph:P002:begin
\section*{A. Vowels and length}\BookHeading{nav-pronunciation-key-0001}{subsection}{A. Vowels and length}{0}
% mdparagraph:P002:end

% answer: PRON-P01a
% mdparagraph:P003:begin
\textbf{A1. for, “for”; fōr, “went”.} \textbf{For} has the short rounded \textbf{o} described in §2. \textbf{Fōr} has long \textbf{ō}: hold the steady initial vowel of \emph{go}, without its final glide. The macron changes duration; it does not mark stress or authorize copying the vowel of the English gloss.
% mdparagraph:P003:end

% answer: PRON-P01b
% mdparagraph:P004:begin
\textbf{A2. węnde, “went”; wēnde, “thought; supposed”.} Short \textbf{ę} has the vowel of \emph{bet}. Long \textbf{ē} has the steady initial vowel of \emph{day}, held without the glide. The hook adds no separate sound. Both words have two syllables, with an audible light \textbf{e} in the second.
% mdparagraph:P004:end

% answer: PRON-P01c
% mdparagraph:P005:begin
\textbf{A3. dǣd, “deed”; wīd, “wide”; hȳdan, “hide; conceal”.} All three marked vowels are long. \textbf{Ǣ} has the held vowel of \emph{cat}; \textbf{ī} has the held vowel of \emph{see}. For \textbf{ȳ}, keep that same high front tongue position but round the lips as for \textbf{oo} in \emph{food}. Moving the tongue back would turn it into ordinary \textbf{ū}. \textbf{Dǣd} and \textbf{wīd} have one syllable each; \textbf{hȳdan} has two. Do not substitute the diphthongs of English \emph{wide} or \emph{hide}.
% mdparagraph:P005:end

% answer: PRON-P01d
% mdparagraph:P006:begin
\textbf{A4.} Each diphthong is one glide within one syllable. The macron lengthens its first part.
% mdparagraph:P006:end

% mdtable:T01:begin
% course-table-style: 2026-10
% course-panel:T01-P01:begin
\par\begingroup
\fontsize{10.5}{13.2}\selectfont
\setlength{\LTleft}{0pt}\setlength{\LTright}{\fill}
\noindent
\begin{tabular}{@{}>{\raggedright\arraybackslash}p{34.265mm}>{\raggedright\arraybackslash}p{86.923mm}@{}}
\toprule
\textbf{Word and supplied meaning} & \textbf{Diphthong and syllables} \\
\midrule
eall, all & Short \textbf{ea}: \emph{cat} vowel toward \textbf{oo}; one syllable. \\
scēap, sheep & Long \textbf{ēa}: held \emph{cat} vowel toward \textbf{oo}; one syllable. \\
heorte, heart & Short \textbf{eo}: steady \emph{day}-type vowel toward \textbf{oo}; two syllables, including final \textbf{e}. \\
frēond, friend & Long \textbf{ēo}: held steady \emph{day}-type vowel toward \textbf{oo}; one syllable. \\
scield, shield & Short \textbf{ie}: \emph{see} vowel toward \textbf{oo}; one syllable. \\
gīet, yet; still & Long \textbf{īe}: held \emph{see} vowel toward \textbf{oo}; one syllable. \\
\bottomrule
\end{tabular}
\par\endgroup
% course-panel:T01-P01:end
% mdtable:T01:end

% mdparagraph:P007:begin
Neither two vowel letters nor extra duration requires an extra syllable. Final \textbf{ll} in \textbf{eall} has ordinary single-consonant duration under the guide's convention.
% mdparagraph:P007:end

\reserve{4\baselineskip}
% mdparagraph:P008:begin
\section*{B. Consonants and unfamiliar sounds}\BookHeading{nav-pronunciation-key-0002}{subsection}{B. Consonants and unfamiliar sounds}{0}
% mdparagraph:P008:end

% answer: PRON-P02a
% mdparagraph:P009:begin
\textbf{B1. cēosað, “choose”.} Initial \textbf{c} is \textbf{ch}, as in \emph{chin}, before the front diphthong \textbf{ēo}. This vetted example follows the ordinary initial pattern. The following diphthong remains one long glide.
% mdparagraph:P009:end

% answer: PRON-P02b
% mdparagraph:P010:begin
\textbf{B2. micel, “large; great”; micles, a declined form.} In \textbf{micel}, \textbf{c} is \textbf{ch}: it follows \textbf{i} and precedes \textbf{e}. In \textbf{micles}, it is \textbf{k}: it immediately precedes the consonant \textbf{l}. The same word can change its consonant sound when the form changes. Merely noticing the preceding \textbf{i} misses the second condition.
% mdparagraph:P010:end

% answer: PRON-P02c
% mdparagraph:P011:begin
\textbf{B3. gān, “go; walk”; glæd, “glad”.} Both begin with stopped \textbf{g}, as in \emph{go}. In \textbf{gān}, the following \textbf{ā} is a back vowel. In \textbf{glæd}, the immediate neighbour is \textbf{l}, so the later \textbf{æ} does not put initial \textbf{g} immediately before a front vowel. Neither enters the \textbf{y} branch; initial position then selects the stopped sound.
% mdparagraph:P011:end

% answer: PRON-P02d
% mdparagraph:P012:begin
\textbf{B4. fæger, “beautiful; fair”; regn, “rain”.} Both have \textbf{g} pronounced as \textbf{y} in \emph{yes}. It follows a front vowel, \textbf{æ} or \textbf{e}, and no back vowel immediately follows it. In \textbf{regn}, \textbf{g} comes before \textbf{n}, not after it; the stopped-\textbf{g} condition “after \textbf{n}” therefore does not apply. Do not transfer the \textbf{c}-before-consonant rule to \textbf{g}.
% mdparagraph:P012:end

% answer: PRON-P02e
% mdparagraph:P013:begin
\textbf{B5. āgan, “have; possess”.} Its single medial \textbf{g} lies between back vowels, so it does not enter the \textbf{y} branch. It is neither initial nor after \textbf{n} nor doubled: use continuous voiced \textbf{g}. Bring the back of the tongue near the hard-\textbf{g} position, leave a narrow gap, and let air pass with throat vibration. A complete closure gives the wrong stopped sound; omitting the sound loses a consonant.
% mdparagraph:P013:end

% answer: PRON-P02f
% mdparagraph:P014:begin
\textbf{B6. brycg, “bridge”; singan, “sing”.} \textbf{Cg} has the final consonant sound of \emph{bridge}. \textbf{Ng} has the sequence heard in \emph{finger}, including the stopped \textbf{g}; do not reduce it to the end of modern English \emph{sing}. The short \textbf{y} of \textbf{brycg} uses the tongue position of \emph{bit} with rounded lips. Combination rules take priority over single-letter rules.
% mdparagraph:P014:end

% answer: PRON-P02g
% mdparagraph:P015:begin
\textbf{B7. riht, “right; direct”; dohtor, “daughter”.} After front \textbf{i} in \textbf{riht}, raise the middle of the tongue toward the roof of the mouth for the front, unvoiced continuous sound. After back \textbf{o} in \textbf{dohtor}, use the back of the tongue near the \textbf{k} position. Keep a gap for breath in both; neither is an ordinary \textbf{k} stop. Both following \textbf{t} sounds are audible. \textbf{Dohtor} has two syllables and a short \textbf{o}.
% mdparagraph:P015:end

% answer: PRON-P02h
% mdparagraph:P016:begin
\textbf{B8. scūfan, “thrust”; cwæð, “said”.} \textbf{Sc} gives the single \textbf{sh} sound of \emph{ship}, even before \textbf{ū}. \textbf{Cw} gives \textbf{k} followed by \textbf{w}, joined without an inserted vowel. Applying a single-\textbf{c} rule inside \textbf{sc} would wrongly produce a separate \textbf{k}.
% mdparagraph:P016:end

\reserve{4\baselineskip}
% mdparagraph:P017:begin
\section*{C. Voice}\BookHeading{nav-pronunciation-key-0003}{subsection}{C. Voice}{0}
% mdparagraph:P017:end

% answer: PRON-P03a
% mdparagraph:P018:begin
\textbf{C1. ǣfre, “ever; always”.} \textbf{F} is voiced, as in \emph{van}. Its neighbours are voiced \textbf{ǣ} and \textbf{r}. The rule requires voiced sounds on both sides, not two vowels.
% mdparagraph:P018:end

% answer: PRON-P03b
% mdparagraph:P019:begin
\textbf{C2. fīf, “five”.} Both \textbf{f} sounds are unvoiced, as in \emph{fan}: one is initial and the other final. Modern English \emph{five} does not license a final \textbf{v} here.
% mdparagraph:P019:end

% answer: PRON-P03c
% mdparagraph:P020:begin
\textbf{C3. cēosað, “choose”.} Medial \textbf{s} is voiced, as in \emph{zip}, between the voiced \textbf{ēo} glide and \textbf{a}. Final \textbf{ð} is unvoiced, as in \emph{thin}. Decide each consonant from its own environment.
% mdparagraph:P020:end

% answer: PRON-P03d
% mdparagraph:P021:begin
\textbf{C4. hæfð, “has”.} Final \textbf{ð} is unvoiced, as in \emph{thin}. The preceding \textbf{f} is next to that voiceless consonant, so it is unvoiced too, as in \emph{fan}. The preceding vowel alone does not make \textbf{f} voiced.
% mdparagraph:P021:end

% answer: PRON-P03e
% mdparagraph:P022:begin
\textbf{C5. cwæð, “said”.} Final \textbf{ð} is unvoiced, as in \emph{thin}. Writing thorn instead would not change the sound: \textbf{þ} and \textbf{ð} do not encode a voiced/\allowbreak{}unvoiced distinction. This question compares spellings; it does not authorize changing the course's edited word.
% mdparagraph:P022:end

% answer: PRON-P03f
% mdparagraph:P023:begin
\textbf{C6. hǣlðu, “health”; ǣfre, “ever; always”.} The supplied exception gives \textbf{ð} in \textbf{hǣlðu} unvoiced \textbf{th}, although \textbf{l} and \textbf{u} are voiced. In \textbf{ǣfre}, the ordinary rule gives voiced \textbf{f} between \textbf{ǣ} and \textbf{r}. A word-specific pronunciation takes priority over the default. Recognizing that limit is the task; deriving the exceptional sound from visible spelling is not required.
% mdparagraph:P023:end

\reserve{4\baselineskip}
% mdparagraph:P024:begin
\section*{D. Syllables and stress}\BookHeading{nav-pronunciation-key-0004}{subsection}{D. Syllables and stress}{0}
% mdparagraph:P024:end

% answer: PRON-P04a
% mdparagraph:P025:begin
\textbf{D1. sittan, “sit”.} Two syllables, with stress on the first. Hold the \textbf{t} closure longer before releasing into the second syllable. \textbf{I} stays short; holding \textbf{tt} neither lengthens it nor adds a vowel or syllable.
% mdparagraph:P025:end

% answer: PRON-P04b
% mdparagraph:P026:begin
\textbf{D2. sunne, “sun”.} Two syllables, with stress on the first. Hold medial \textbf{nn} longer, as with the \textbf{n} across \emph{pen-knife}. Its duration adds no syllable. Keep the final light \textbf{e} audible; modern English's one-syllable \emph{sun} is not the model for the whole word.
% mdparagraph:P026:end

% answer: PRON-P04c
% mdparagraph:P027:begin
\textbf{D3. gewrit, “letter; written document”.} Two syllables; stress \textbf{writ}, after unstressed \textbf{ge-}. Pronounce \textbf{w+r} without an inserted vowel. There is no doubled consonant. The vowel \textbf{i} remains short despite receiving the stress.
% mdparagraph:P027:end

% answer: PRON-P04d
% mdparagraph:P028:begin
\textbf{D4. forlǣtan, “abandon; forsake”.} Three syllables, with stress on \textbf{lǣ} in the root. The supplied division identifies unstressed \textbf{for-}. The macron says that \textbf{ǣ} is long; the prefix/\allowbreak{}root rule determines stress. There is no doubled consonant, and final \textbf{-tan} remains audible.
% mdparagraph:P028:end

\reserve{4\baselineskip}
% mdparagraph:P029:begin
\section*{E. A sentence adapted from Cook}\BookHeading{nav-pronunciation-key-0005}{subsection}{E. A sentence adapted from Cook}{0}
% mdparagraph:P029:end

% mdparagraph:P030:begin
The supplied meaning is “God created heaven and earth.” These questions apply sound rules; they do not test independent translation.
% mdparagraph:P030:end

% answer: PRON-P05a
% mdparagraph:P031:begin
\textbf{E1. gescēop, “created”.} Initial \textbf{g} before \textbf{e}, in the supplied prefix \textbf{ge-}, has the \textbf{y} sound of \emph{yes}. The combination \textbf{sc} has \textbf{sh}. There are two syllables: the light prefix and stressed \textbf{scēop}, whose \textbf{ēo} is one long glide. Keep the prefix vowel; \textbf{g} and \textbf{sc} are not adjacent consonant sounds.
% mdparagraph:P031:end

% answer: PRON-P05b
% mdparagraph:P032:begin
\textbf{E2. heofonan, “heaven”.} \textbf{F} is voiced, as in \emph{van}, between the voiced \textbf{eo} glide and \textbf{o}. There are three syllables: the first contains \textbf{eo}, and the next two contain \textbf{o} and \textbf{a}. The glide does not take two syllables; neither unstressed vowel is silent.
% mdparagraph:P032:end

% answer: PRON-P05c
% mdparagraph:P033:begin
\textbf{E3. eorðan, “earth”.} \textbf{Ð} is voiced, as in \emph{this}, between voiced \textbf{r} and \textbf{a}. The rule is “between voiced sounds”, so a voiced consonant is as relevant as a vowel. Keep the initial short \textbf{eo} glide and the audible second syllable.
% mdparagraph:P033:end


% book-source-body:end
\par\label{end-pronunciation-key}
\clearpage\endgroup
